\subsection{Diagonalization of compact endomorphisms}

THEOREM 1. --- Suppose K = C. Let u be a compact and normal endomorphism of E. There exists an orthonormal basis B of E such that u is diagonal in the basis B.

The set \(\mathrm{Sp}_{\mathrm{s}}(u)\) is countable, and it does not contain 0 if E is infinite-dimensional (III, p. 90, prop. 5, \(b\) ). For every element \(\lambda \in \mathrm{Sp}_{\mathrm{s}}(u)\) , denote by \(\mathrm{N}_{\lambda}\) the nilspace of \(u - \lambda 1_{\mathrm{E}}\) . It is finite-dimensional (loc. cit.) and, since \(u\) is normal, it coincides with the eigenspace of \(u\) relative to \(\lambda\) (EVT, V, p. 43, cor. of prop. 8). The spaces \(\mathrm{N}_{\lambda}\) are pairwise orthogonal (I, p. 132, n° 5).

Let F be the subspace of E that is the Hilbert sum of the spaces \(N_{\lambda}\) for \(\lambda \in \mathrm{Sp}_{\mathrm{s}}(u) - \{0\}\) . It is a space of countable type, stable under u since each subspace \(N_{\lambda}\) is stable under u. Since \(N_{\lambda}\) is also the eigenspace of \(u^{*}\) relative to \(\overline{\lambda}\) (EVT, V, loc. cit.), the endomorphism u induces an endomorphism \(\widetilde{u}\) of \(F^{\circ}\) by passing to subspaces. The endomorphism \(\widetilde{u}\) is compact (prop. 3 of III, p. 5) and normal (lemma 4 of I, p. 135). By construction, the sensitive spectrum of \(\widetilde{u}\) is contained in \(\{0\}\) (indeed, every eigenvector for \(\widetilde{u}\) would be one for u, hence would belong to one of the spaces \(N_{\lambda}\) if the corresponding eigenvalue were nonzero). Thus the spectral radius of \(\widetilde{u}\) is zero, whence \(\widetilde{u}=0\) since \(\widetilde{u}\) is normal (cor. 1 of I, p. 108). We therefore have \(F^{\circ}\subset\operatorname{Ker}(u)\) .

For every \(\lambda \in \mathrm{Sp}_{\mathrm{s}}(u)\) , let \(B_{\lambda}\) be an orthonormal basis of \(N_{\lambda}\) and let \((e_j)_{j \in J}\) be the family formed by the union of the \(B_{\lambda}\) ; it is an orthonormal basis of F. Let B be the union of \((e_j)_{j \in J}\) and of an orthonormal basis of \(F^{\circ}\) . This is an orthonormal basis of E and \(u\) is diagonal in the basis B.

COROLLARY 1. --- Let u be a compact Hermitian endomorphism of E. There exists an orthonormal basis B of E such that u is diagonal in the basis B and its eigenvalues are real.

If K = C, this follows immediately from the theorem. Suppose K = R. The space E is an R-structure on \(\mathrm{E}_{(\mathbf{C})}\) (cf. A, II, p. 119). The endomorphism \(u_{(\mathbf{C})}\) of \(\mathrm{E}_{(\mathbf{C})}\) is compact (remark 4 of III, p. 2) and Hermitian. Let \(B = (e_j)_{j \in J}\) be an orthonormal basis of \(\mathrm{E}_{(\mathbf{C})}\) such that \(u_{(\mathbf{C})} \in \mathcal{D}_{\mathrm{B}}(\mathrm{E}_{(\mathbf{C})})\) and let \(\lambda\) be the family of eigenvalues of \(u_{(\mathbf{C})}\) (theorem 1). We have \(\lambda \in R^{J}\) (I, p. 106, prop. 4); since the linear map \(u_{(\mathbf{C})}\) is R-rational, the eigenspace of \(u_{(\mathbf{C})}\) relative to \(\lambda_j\) is R-rational for every \(j \in J\) (cf. A, V, p. 60, prop. 6). Hence there exists a basis of it belonging to E, and a fortiori, there exists an orthonormal basis in E. The union of these bases is an orthonormal basis \(B_R\) of E such that \(u \in \mathcal{D}_{\mathrm{B}_R}(\mathrm{E})\) .

COROLLARY 2. --- Let F be a Hilbert space and let u be a compact continuous linear map from E to F. There exists a countable set I, an orthonormal basis \((e_{i})_{i\in\mathrm{I}}\) of the initial space \(\operatorname{Ker}(u)^{\circ}\) of u, an orthonormal family \((f_{i})_{i\in\mathrm{I}}\) of F, and a family \((\alpha_{i})_{i\in\mathrm{I}}\in(\mathbf{R}_{+}^{*})^{\mathrm{I}}\) such that \(u(e_{i})=\alpha_{i}f_{i}\) for all \(i\in I\) .

Let \(v = u^{*} \circ u\) . This is a compact, positive, hence Hermitian, endomorphism of E. By corollary 1, there exists an orthonormal basis \((e_{j})_{j \in \mathrm{J}}\) of E such that \(v\) is diagonal in this basis. The family \((\lambda_{j})_{j \in \mathrm{J}}\) of its eigenvalues is contained in \(\mathbf{R}_{+}^{\mathrm{J}}\) . Set \(\alpha_{j} = \sqrt{\lambda_{j}}\) for all \(j \in \mathrm{J}\) . Let I be the set of \(j \in \mathrm{J}\) such that \(\alpha_{j} \neq 0\) . This is a countable set since \(v\) is compact. The family \((e_{i})_{i \in \mathrm{I}}\) is an orthonormal basis of the initial space of \(v\) which is the initial space \(\operatorname{Ker}(u)^{\circ}\) of \(u\) (EVT, V, p. 43, prop. 8). Set \(f_{i} = \frac{1}{\alpha_{i}} u(e_{i})\) for \(i\in \mathrm{I}\) . For any \(i\) and \(j\) in I, we have

\[
\langle f _ {i} | f _ {j} \rangle = \frac {1}{\alpha_ {i} \alpha_ {j}} \langle u (e _ {i}) | u (e _ {j}) \rangle = \frac {1}{\alpha_ {i} \alpha_ {j}} \langle v (e _ {i}) | e _ {j} \rangle = \frac {\lambda_ {i}}{\alpha_ {i} \alpha_ {j}} \langle e _ {i} | e _ {j} \rangle ,
\]

whence it follows that the family \((f_{i})_{i\in\mathrm{I}}\) is orthonormal in F. The corollary follows, since \(u(e_{i})=\alpha_{i}f_{i}\) for all \(i\in I\) .

DEFINITION 2. --- With the notations of the corollary, the family \((\alpha_{i})_{i\in I}\) is the family of singular values of \(u\) relative to the orthonormal basis \((e_{i})_{i\in I}\) of the initial space of \(u\) .

Remarks. --- 1) This corollary generalizes th. 2 of EVT, V, p. 54, which corresponds to Hilbert--Schmidt maps.

\begin{enumerate}
\def\labelenumi{\arabic{enumi})}
\setcounter{enumi}{1}
\tightlist
\item
  With the notations of the corollary, we have the formula
\end{enumerate}

\[
u (x) = \sum_ {i \in I} \alpha_ {i} \langle e _ {i} | x \rangle f _ {i}\tag{1}
\]

for all \(x \in E\) .

\begin{enumerate}
\def\labelenumi{\arabic{enumi})}
\setcounter{enumi}{2}
\tightlist
\item
  Let \(u\) be a compact positive endomorphism of E. Let \(B = (f_i)_{i \in I}\) be an orthonormal basis of E such that \(u\) is diagonal in the basis B (theorem 1), and let \((\lambda_i)_{i \in I}\) be the family of eigenvalues of \(u\) in this basis. Let J be the set of \(i \in I\) such that \(\lambda_i > 0\) ; the family \((e_i)_{i \in J}\) is an orthonormal basis of the space \(\text{Ker}(u)^{\circ}\) . For every \(i \in J\) , set \(e_i = f_i\) and \(\alpha_i = \lambda_i\) . We obtain \(u(e_i) = \alpha_if_i\) : the family \((\alpha_i)_{i \in J}\) is the family of singular values of \(u\) relative to the basis \((e_i)_{i \in J}\) .
\end{enumerate}

